% ============================================================================
% LANGENTWURF LE-004: Arbeitsumgebung iPhone
% Profil: S (Senioren 65+) | Tier: 2 (Standard)
% ============================================================================

\documentclass[11pt,a4paper]{article}
\usepackage[utf8]{inputenc}
\usepackage[T1]{fontenc}
\usepackage[ngerman]{babel}
\usepackage{geometry}
\usepackage{fancyhdr}
\usepackage{lastpage}
\usepackage{xcolor}
\usepackage{tcolorbox}
\usepackage{enumitem}
\usepackage{longtable}
\usepackage{hyperref}
\usepackage{amssymb}

\geometry{left=2.5cm, right=2.5cm, top=2.5cm, bottom=2.5cm}
\definecolor{fach}{HTML}{b35c00}
\definecolor{gray}{HTML}{666666}
\definecolor{lightgray}{HTML}{f5f5f5}
\definecolor{successgreen}{HTML}{2a7a4b}

\tcbuselibrary{skins,breakable}
\newtcolorbox{scorebox}{ colback=white, colframe=fach, fonttitle=\bfseries, title=Didaktik-Score, breakable }

\pagestyle{fancy}
\fancyhf{}
\fancyhead[L]{\small\textcolor{gray}{Digitale Grundlagen -- 004 Arbeitsumgebung iPhone}}
\fancyhead[R]{\small\textcolor{gray}{LE Tier 2 / Profil S}}
\fancyfoot[C]{\small\thepage{} / \pageref{LastPage}}
\renewcommand{\headrulewidth}{0.4pt}

\begin{document}

\begin{titlepage}
    \centering
    \vspace*{2cm}
    {\Large\textcolor{gray}{Digitale Grundlagen -- Senioren 65+}}\\[2cm]
    {\Huge\textcolor{fach}{\textbf{Langentwurf}}}\\[0.5cm]
    {\large Tier 2 | Profil S}\\[2cm]
    {\LARGE\textbf{004 -- Arbeitsumgebung iPhone}}\\[0.5cm]
    {\large Modul A: Geräte \& Zugänge}\\[2cm]
    \begin{tabular}{rl}
        \textbf{Zeitrahmen:} & 60--90 Minuten \\
        \textbf{Vorkenntnisse:} & iPhone einschalten, entsperren \\
    \end{tabular}
    \vfill
    {\normalsize Stand: \today}
\end{titlepage}

\tableofcontents
\newpage

\section{Bedingungsanalyse}

\subsection{Ausgangssituation}
\begin{itemize}
    \item Teilnehmer besitzen iPhone, nutzen es aber nur für Telefonate
    \item Home-Bildschirm ist ``Chaos'' oder unverändert seit Einrichtung
    \item Gesten sind unbekannt oder unsicher
    \item Kontrollzentrum und Suche werden nicht genutzt
\end{itemize}

\subsection{Typische Aussagen}
\begin{itemize}
    \item ``Ich finde die App nicht mehr''
    \item ``Wie komme ich zurück zum Anfang?''
    \item ``Was sind diese Punkte unten?''
    \item ``Ich wische aus Versehen irgendwo hin''
\end{itemize}

\section{Sachanalyse}

\subsection{Kernkonzepte}
\begin{enumerate}
    \item \textbf{Home-Bildschirm:} Der ``Schreibtisch'' -- Apps, Dock, Seiten
    \item \textbf{Apps:} ``Werkzeuge'' für verschiedene Aufgaben
    \item \textbf{Gesten:} Tippen, Wischen, Halten, Zoomen
    \item \textbf{Kontrollzentrum:} Schnellzugriff auf WLAN, Helligkeit etc.
    \item \textbf{Einstellungen:} Die ``Schaltzentrale''
    \item \textbf{Suche:} Alles finden (Apps, Kontakte, Einstellungen)
\end{enumerate}

\subsection{Alltagsanalogien}
\begin{tabular}{|l|l|}
\hline
Home-Bildschirm & Schreibtisch \\
Apps & Werkzeuge \\
Dock & Die 4 wichtigsten Werkzeuge immer griffbereit \\
Gesten & Handgriffe \\
Kontrollzentrum & Schalttafel im Auto \\
Einstellungen & Sicherungskasten \\
\hline
\end{tabular}

\section{Didaktische Analyse}

\subsection{Stellung in der Reihe}
Einheit 004 ist die erste ``Bedienungs-Einheit''. Nach den konzeptuellen Grundlagen (001-003) geht es jetzt um praktische Handhabung.

\subsection{Legitimation}
\begin{itemize}
    \item Grundvoraussetzung für alle weiteren Einheiten
    \item Reduziert Frustration und Angst vor ``falschem Wischen''
    \item Ermöglicht Selbstständigkeit im Alltag
\end{itemize}

\section{Lernziele}

\begin{tabular}{|c|p{7cm}|c|c|}
\hline
\textbf{Nr.} & \textbf{Die Teilnehmer können...} & \textbf{Stufe} & \textbf{Niveau} \\
\hline
FZ1 & die Elemente des Home-Bildschirms benennen & Wissen & Basis \\
FZ2 & eine App durch Tippen öffnen & Anwenden & Basis \\
FZ3 & zum Home-Bildschirm zurückkehren & Anwenden & Basis \\
FZ4 & das Kontrollzentrum öffnen und schließen & Anwenden & Basis \\
FZ5 & die Suche nutzen, um Apps zu finden & Anwenden & Basis \\
FZ6 & die Schriftgröße in den Einstellungen ändern & Anwenden & Vertiefung \\
\hline
\end{tabular}

\section{Methodische Analyse}

\subsection{Besonderheit: Gesten lernen}
Gesten müssen \textbf{motorisch geübt} werden. Nicht nur erklären, sondern:
\begin{itemize}
    \item Vormachen (am eigenen Gerät demonstrieren)
    \item Gemeinsam machen (``Jetzt alle zusammen'')
    \item Selbst üben (mit Checkliste abhaken)
    \item Wiederholen (mehrfach in der Stunde)
\end{itemize}

\subsection{Material}
\begin{itemize}
    \item Seite 1: Home-Bildschirm, Apps, Gesten (Drei-Spalten-Infoblatt)
    \item Seite 2: Kontrollzentrum, Einstellungen, Suche
    \item Seite 3: Praktische Übungen + Checkliste
\end{itemize}

\section{Verlaufsplanung}

\begin{longtable}{|p{1cm}|p{2cm}|p{5cm}|p{3cm}|p{2cm}|}
\hline
\textbf{Zeit} & \textbf{Phase} & \textbf{Inhalt} & \textbf{TN-Aktivität} & \textbf{Material} \\
\hline
0--10 & Einstieg & ``Was sehen Sie, wenn Sie Ihr iPhone entsperren?'' Begriffe sammeln. & Beschreiben & -- \\
\hline
10--25 & Home-Bildschirm & Elemente erklären, Analogie Schreibtisch, Seiten wischen üben & Mitmachen & S.1 \\
\hline
25--40 & Gesten & Die 6 Grundgesten einzeln üben, jede mehrfach & Üben am Gerät & S.1 \\
\hline
40--55 & Kontrollzentrum & Öffnen/Schließen, wichtigste Schalter, Taschenlampe testen & Ausprobieren & S.2 \\
\hline
55--70 & Suche \& Einstellungen & Suche öffnen, App suchen, Schriftgröße ändern & Selbst durchführen & S.2 \\
\hline
70--80 & Übungen & Gesten-Checkliste durcharbeiten & Abhaken & S.3 \\
\hline
\end{longtable}

\section{Lösungen}

\textbf{Teil A: Begriffe zuordnen}\\
1-D, 2-C, 3-A, 4-B, 5-E

\textbf{Teil C: Was machst du?}
\begin{enumerate}
    \item WLAN schnell ein/aus: Kontrollzentrum öffnen (von oben rechts wischen)
    \item App nicht gefunden: Suche öffnen (von Mitte nach unten wischen) oder App-Mediathek
    \item Schrift zu klein: Einstellungen $\rightarrow$ Anzeige \& Helligkeit $\rightarrow$ Textgröße
\end{enumerate}

\section{Didaktik-Score}

\begin{scorebox}
\begin{tabular}{|c|l|c|c|p{3.5cm}|}
\hline
\# & Dimension & Gew. & Score & Notiz \\
\hline
1 & Differenzierung & 15\% & 8/10 & Basis-Gesten vs. Vertiefung \\
2 & Taxonomie & 10\% & 10/10 & Anwenden dominiert (richtig!) \\
3 & Alltagsrelevanz & 10\% & 10/10 & Fundamental für alles \\
4 & Aktivierung & 10\% & 10/10 & Ständiges Üben am Gerät \\
5 & Scaffolding & 10\% & 9/10 & Schrittweise, Analogien \\
6 & Feedback & 10\% & 9/10 & Checkliste zum Abhaken \\
7 & Sprachliche Klarheit & 10\% & 10/10 & Sehr einfach, Analogien \\
8 & Motivation & 10\% & 9/10 & Erfolgserlebnisse durch Üben \\
9 & Barrierefreiheit & 10\% & 9/10 & Große Darstellung, Farben \\
10 & Praktikabilität & 5\% & 8/10 & 80 Min evtl. lang \\
\hline
\multicolumn{3}{|r|}{\textbf{GESAMT:}} & \textbf{9.25/10} & \textcolor{successgreen}{\textbf{FREIGABE}} \\
\hline
\end{tabular}
\end{scorebox}

\end{document}
